\begin{proof}[Proof of Theorem~\ref{thm:kernel}]
By Lemma~\ref{lem:scope} both systems reduce over the same set of series, so it suffices
to match the two reductions on one series $s$ and then the space aggregation.

For the time reduction, (1) makes $[t,t+\Delta)$ and the PromQL window at
$u(t)=t+\Delta-\eta$ cover the same samples, by the index argument of
Proposition~\ref{prop:offset}; (2) makes it the same width. Hypothesis~(3) fixes the metric's
type. The lowering writes $\theta=\text{counter}$ where the query carries the evidence
\S\ref{sec:lowering} names and $\theta=\text{unknown}$ otherwise --- it never writes
$\text{gauge}$ --- and the emitter routes $\text{unknown}$ to the gauge rule. So the two
cases below are the emitter's two rules rather than three values of $\theta$, and the
case that carries almost the whole corpus is the first of them reached through
$\text{unknown}$, which is also why \texttt{METRIC\_TYPE\_UNKNOWN} marks every
translation.

If $\theta$ is a gauge, $\phi_{\mathsf{zero}}(s)=s$, so Datadog's stored points in the
bucket are exactly the samples $s|_{[t,t+\Delta)}$ that
\texttt{$\mathit{op}_t$\_over\_time} reduces at that instant, and the two reductions are
the same commutative operation on the same multiset.

If $\theta$ is a counter and $\mathit{op}_t=\mathsf{sum}$, Datadog reports
$\Sigma_\Delta(t)$ and the emitted form equals it by Proposition~\ref{prop:offset}
under (4)'s first branch and (5), or by Proposition~\ref{prop:anchored} under (4)'s
second, which needs neither --- but which needs a wider lookback in their place, $u-t_a$
for $t_a$ the anchor rather than the two short distances of
Proposition~\ref{prop:offset}. An earlier version of this line said only ``which needs
neither'', and the theorem's tail asked for Proposition~\ref{prop:offset}'s distances
whichever branch was taken, so the anchored route was being granted a hypothesis
narrower than the one it uses. The discussion after the theorem had the fact right while
the hypothesis list did not.

If $\theta$ is a counter and $\mathit{op}_t=\mathsf{avg}$, Datadog reports the mean of
the deltas it stores in the bucket. The emitted form is
$\Sigma_\Delta(t)/\texttt{count\_over\_time}$, whose numerator is that sum by the
previous case and whose denominator counts the samples in the window. By (5) the series
began before $t$, so every sample in $[t,t+\Delta)$ contributes exactly one delta and
the two counts agree; the quotient is the mean.

If $\theta$ is a counter and $\mathit{op}_t=\mathsf{count}$, Datadog reports how many
deltas it stores in the bucket and the emitted \texttt{count\_over\_time} reports how
many samples the window holds; (5) makes those the same number, by the same
correspondence. For $\mathit{op}_t\in\{\mathsf{min},\mathsf{max}\}$ over a counter no
rule exists --- the individual deltas are reachable only through a subquery over a
sample grid the query does not state --- so the emitter refuses,
$q\notin\mathrm{dom}(T)$, and there is nothing to prove.
The $\mathit{per\_second}$ divisor is $\Delta$ on both sides.

For the space reduction, Definition~\ref{def:ir-sem}'s $\mathrm{Agg}$ is what PromQL's
$\mathit{op}_s$ \texttt{by} $(\overline{k})$ computes and what Datadog's
$\mathit{op}_s{:}$ with \texttt{by \{$\overline{k}$\}} computes, on inputs just shown
equal. This is the step the grouping-key injectivity hypothesis is for, and the step at
which it was never invoked: two of $q$'s grouping keys that mangle to one PromQL label
name would merge two of Datadog's groups into one of PromQL's, so the two aggregations
would not compute the same function however the inputs agree. The hypothesis is needed
and Remark~\ref{rem:mangling} shows it is not vacuous --- \texttt{cluster-name} against
\texttt{cluster\_name} is a measured instance --- so the gap was in the discharge and not
in the statement.

One hypothesis is stronger than this proof requires, in the other direction. Hypothesis
(5) is stated for every form that reads a counter, and the sum case under (4)'s second
branch does not use it: Proposition~\ref{prop:anchored} covers series with no sample
before the bucket, which is exactly what (5) excludes. \S\ref{sec:emission} says the
anchored form ``avoids'' \texttt{PHI\_SERIES\_START} for that reason, so the hypothesis
list and the narrative disagree, and the list is the cautious one --- it asks for more
than is needed, which costs a reader a condition rather than costing the theorem its
truth. It is left as stated and flagged here, because the emitter records the code on
those forms whether or not this proof consumes it, and $C$ is defined as what the
emitter records.

That $C$ contains nothing further is by inspection of the emitter. The inspection is
performed rather than asserted, by
\texttt{tests\slash test\_paper\_formulas.py}: the seven
codes the hypotheses above name --- six hypotheses, the fourth naming two --- are the
only ones it records anywhere.
\end{proof}
